\adnote{This chapter description from the proposal will function as the "roadmap" for the structure and contents of this dissertation chapter.}

%========================
% Chapter 1 — Introduction (Proposal)
%========================

% \chapter{Introduction} % uncomment if your class uses \chapter

This chapter establishes the problem, rationale, and proposed contributions. It frames digital tools as culturally loaded and non-neutral and situates audio software within Eurocentric defaults that shape perception, interaction, and creative possibility. It introduces \textit{Speculocultural Technopoiesis} (ST) as the guiding framework, articulates the central research question and sub-questions, and motivates a shift toward culturally grounded tool modification where Black sonic epistemologies acts as and example for becoming computational logic. Finally, it previews the methodology and case studies and explains the significance for digital media, HCI/interaction design, and Black studies as related to the study.

\subsubsection{Objectives}
\begin{itemize}[nosep]
	\item Introduce ST and state the central research question.
	\item Establish the problem (non-neutral tools) and articulate the significance.
	\item Preview methodology (autoethnography, collaborative sessions, research-creation, critical technical practice).
	\item Identify scope, assumptions, and biases.
	\item Preview practice-based case studies and their role in answering the research question(s).
\end{itemize}

\subsubsection{Intended Sections/Subsections}
\begin{itemize}[nosep]
	\item Background and Context
	\item Problem Statement (Unsettling Foundations)
	\item Purpose and Research Questions
		\begin{itemize}[nosep]
			\item Primary Research Question
			\item Guiding Sub-Questions
		\end{itemize}
	\item Guiding Framework: Sonic Speculocultural Technopoiesis
		\begin{itemize}[nosep]
			\item \textbf{Unsettling} — surface friction between cultural knowledge and tool assumptions; identify sites for change.
			\item \textbf{Remixing} — materially modify tools/workflows; encode cultural relations (e.g., polyrhythmic time, call–response).
			\item \textbf{Embodying} — test for cultural authenticity and experiential fit via practice; validate what the modification enables.
			\item \textbf{Transmitting} — document, name, and share methods and code so knowledge circulates with its cultural reasoning.
		\end{itemize}
		\begin{itemize}[nosep]
			\item \textbf{Fugitive Speculation} (refusal of intended uses; speculative alternatives),
			\item \textbf{Material Memory} (embedding cultural archives in parameters/structures),
			\item \textbf{Sensory Worlding} (multisensory, movement- and sound-first interaction),
			\item \textbf{Intimate Witnessing} (autoethnography, cohort validation, ethical stance).
		\end{itemize}
	\item Significance and Contributions
	\item Scope, Delimitations, and Assumptions
	\item Methodological Preview
	\item Practice-Based Works and Case Studies
\end{itemize}
%-------------------------------------------------------------------------------
\begin{comment}
\subsubsection{Background and Context}
Across contemporary creative practice, digital tools are often treated as neutral infrastructures. Yet both philosophy of technology and critical race/technology studies show that tools mediate experience and encode cultural assumptions (e.g., default temporal grids, tuning systems, workflows). In audio software, these defaults privilege particular epistemologies of time, harmony, authorship, and interaction. The issue is not merely aesthetic: defaults delimit what creators can easily perceive and do, channeling practice toward certain sonic possibilities while obscuring others.

\subsubsection{Problem Statement (Unsettling Foundations)}
The central problem is that prevailing audio tools normalize Western/Eurocentric logics as if they were universal. These logics become “taken for granted” in the interface, shaping perception, interaction, and pedagogy. For practitioners working from Black sonic traditions—where polyrhythm, call-and-response, groove, and improvisatory assemblage organize practice—these defaults can misfit, constrain, or render important cultural techniques invisible or as alien. The dissertation responds by proposing a method to \textit{unsettle} these defaults and rebuild them from situated cultural logics.

\subsubsection{Purpose and Research Questions}
\textbf{Purpose.} To develop and demonstrate a culturally grounded method—\textit{Speculocultural Technopoiesis}—for modifying or designing audio interaction so that Black sonic epistemologies operate as first-order computational logic.

\vspace{1.0em}
\noindent\textbf{Primary Research Question.}
\begin{quote}
	How can Black sonic epistemologies be operationalized as computational logic in audio tool design so that interaction, perception, and practice are reconfigured at the level of interface and algorithm?
\end{quote}

\noindent\textbf{Guiding Sub-Questions.}
\begin{enumerate}
	\item Which embedded assumptions in current audio tools most directly conflict with or obscure Black sonic practices?
	\item What modifications to interaction patterns, representations, and algorithms make those practices legible and generative in software?
	\item How does collaborative, practice-based making with artist–technologists reveal, test, and refine these modifications?
	\item In what ways does the modified tooling reshape creative outcomes and the theoretical account of mediation?
\end{enumerate}

\subsubsection{Guiding Framework: Speculocultural Technopoiesis (ST)}
ST is a cyclical method linking cultural imagination and embodied making. It proceeds through four movements:

\begin{itemize}
	\item \textbf{Unsettling} — surface friction between cultural knowledge and tool assumptions; identify sites for change.
	\item \textbf{Remixing} — materially modify tools/workflows; encode cultural relations (e.g., polyrhythmic time, call–response).
	\item \textbf{Embodying} — test for cultural authenticity and experiential fit via practice; validate what the modification enables.
	\item \textbf{Transmitting} — document, name, and share methods and code so knowledge circulates with its cultural reasoning.
\end{itemize}

\noindent Four cross-cutting territories guide design decisions:
\begin{itemize}
	\item \textbf{Fugitive Speculation} (refusal of intended uses; speculative alternatives),
	\item \textbf{Material Memory} (embedding cultural archives in parameters/structures),
	\item \textbf{Sensory Worlding} (multisensory, movement- and sound-first interaction),
	\item \textbf{Intimate Witnessing} (autoethnography, cohort validation, ethical stance).
\end{itemize}

\subsubsection{Significance and Contributions}
\textbf{Theoretical.} Re-frames technological mediation as culturally specific; links Black studies and philosophy of technology through a concrete design method. \\
\textbf{Methodological.} Integrates research-creation, critical technical practice, and collaborative (small-cohort) making into a reproducible cycle for tool redesign. \\
\textbf{Practical.} Produces concrete interaction patterns and algorithmic models (e.g., polyrhythmic time bases; call–response interaction graphs) that can inform DAWs, plugins, or bespoke interfaces. \\
\textbf{Field Impact.} Offers HCI/interaction design a template for culturally grounded defaults; contributes to sonic interaction design with procedures rooted in Black epistemologies.

\subsubsection{Scope, Delimitations, and Assumptions}
The project focuses on audio interaction (not full production ecosystems), using targeted prototypes and modifications to expose leverage points. Findings are intentionally situated (small cohort, practice-centered) and prioritize cultural authenticity over generality. Assumptions include the value of autoethnographic insight and the legitimacy of artistic outcomes as research knowledge.

\subsubsection{Methodological Preview}
The study employs (1) autoethnography to document the researcher’s sonic practice; (2) collaborative small-cohort sessions with Black artist–technologists to co-design and test interactions; (3) research-creation to treat making as knowledge production; and (4) critical technical practice to keep theory and implementation in continuous dialogue. Data include reflective logs, session recordings/transcripts, design artifacts, and sonic outputs.

\subsubsection{Practice-Based Works and Case Studies}
Two case studies function as research instruments:
\begin{itemize}
	\item \textbf{Buffalo Soldier} — a spatial sonic computing environment demonstrating the full ST cycle via embodied movement-to-sound logics.
	\item \textbf{Did Sun Ra Fear Cryosleep?} — a temporal/structural study exploring speculative recursion and non-linear, Afrofuturist conceptions of time.
\end{itemize}
Each case iterates Unsettling→Remixing→Embodying→Transmitting and feeds results back into the framework.

\subsubsection{Definitions of Key Terms (Selected)}
\begin{description}
	\item[Speculocultural Technopoiesis (ST)] A method where cultural speculation guides embodied technological making to generate new knowledge and tools.
	\item[Black sonic epistemologies] Ways of knowing organized through Black musical practice (e.g., polyrhythm, groove, call–response, improvisation).
	\item[Mediation] The way tools shape perception/action rather than merely transmit; here treated as culturally specific.
	\item[Fugitive Speculation / Material Memory / Sensory Worlding / Intimate Witnessing] Four design territories used to steer modifications and evaluation.
\end{description}



% --- Chapter 1 Key References (keyword-based filter) ---
\begin{refsection}[references.bib]
	\nocite{*}
	\printbibliography[keyword=ch1,heading=subbibliography,title={Key References}]
\end{refsection}
\end{comment}




